\begin{tikzpicture}

  % Definición del actor
  \umlactor[x=0,y=3, name=usuario_movil]{usuario-móvil}
  \umlactor[x=3.5,y=-6]{usuario-escritorio}
  \umlactor[x=6,y=4.5]{camara}
  \umlactor[x=6,y=-3]{microfono}
  
  % Definición de casos de uso
  \begin{umlsystem}[x=6] {}
    \umlusecase[x=5, y=3, width=2cm]{Tomar Fotografia} %usecase-1
    \umlusecase[x=12, y=3, width=2cm]{Preprocesar Imagen} %usecase-2
    \umlusecase[x=12, y=6, width=2cm]{Guardar Imagen} %usecase-3
    \umlusecase[x=5, y=-3, width=2cm]{Grabar Audio} %usecase-4
    \umlusecase[x=12, y=-3, width=2cm]{Guardar Audio} %usecase-5
    \umlusecase[x=12, y=0, width=2cm]{Preprocesar Audio} %usecase-6
    \umlusecase[x=5, y=-6, width=2cm]{Introducir Diagrama} %usecase-7
    \umlusecase[x=12, y=-6, width=2.2cm]{Descomponer Diagrama} %usecase-8
    \umlusecase[x=5, y=-9, width=2cm]{Escribir nota} %usecase-9
    \umlusecase[x=12, y=-9, width=2cm]{Guardar Diagrama} %usecase-10
    \umlusecase[x=12, y=-12, width=2cm]{Guardar Nota} %usecase-11
  \end{umlsystem}
  
  % Conexiones entre el actor y los casos de uso
  \umlassoc{camara}{usecase-1}
  \umlassoc{microfono}{usecase-4}
  
  \umlassoc{usuario-móvil}{usecase-1}
  \umlassoc{usuario-móvil}{usecase-4}
  \umlVHassoc{usuario-móvil}{usecase-9}

  \umlassoc{usuario-escritorio}{usecase-7}
  \umlassoc{usuario-escritorio}{usecase-9}
  \umlassoc{usuario-escritorio}{usecase-5}
  \umlVHassoc{usuario-escritorio}{usecase-3}

  % extend
  \umlinclude{usecase-1}{usecase-2} 
  \umlextend{usecase-3}{usecase-1} 

  \umlextend{usecase-5}{usecase-4} 
  \umlinclude{usecase-4}{usecase-6}

  \umlinclude{usecase-7}{usecase-8} 
  \umlextend{usecase-10}{usecase-7}

  \umlextend{usecase-11}{usecase-9}

\end{tikzpicture}
